% Vietnamese translation for OI Public Library
% Source-File: IOI中国国家候选队论文/2009/毛杰明/母函数的性质及应用.ppt
\documentclass[11pt]{article}
\usepackage{oipl-vn}

\title{Tính chất và ứng dụng của hàm sinh\\{\large Ghi chú theo slide}}
\author{Mao Jieming\\Trường Ngoại ngữ Nam Kinh\\Email: \texttt{maojm517@163.com}\\Mã bưu chính: 210008}
\date{Bài luận đội tuyển quốc gia năm 2009}

\begin{document}
\maketitle

\origtitle{Properties and applications of generating functions}
\sourcepath{IOI-China-Candidate-Team-Papers/2009/Mao-Jieming/generating-functions-slides.ppt}

\section{Mạch trình bày}

Slide giới thiệu hàm sinh như công cụ mã hóa một dãy số thành một biểu thức
hình thức. Các phép toán trên biểu thức tương ứng với phép cộng, chập, chọn
độc lập hoặc đánh dấu số lượng trong bài toán đếm.

\section{Tính chất cần nhớ}

\begin{itemize}
  \item Cộng hàm sinh tương ứng với tách trường hợp.
  \item Nhân hàm sinh tương ứng với ghép hai lựa chọn độc lập, hệ số của
  \(x^n\) là tổng các cách phân rã chỉ số.
  \item Hàm sinh mũ thuận tiện khi đối tượng có nhãn.
  \item Định lý Pólya kết hợp hàm sinh với nhóm hoán vị để đếm cấu hình theo
  đối xứng.
\end{itemize}

\section{Ứng dụng}

Các ví dụ như \textit{Chocolate}, \textit{Sweet} và \textit{Polygon} minh họa
ba kiểu dùng: lập biểu thức rồi lấy hệ số, biến đổi đại số để có công thức
nhanh, và đưa đối xứng vào bài toán đếm. Khi đọc slide, nên theo dõi đối tượng
nào được đánh dấu bằng biến và hệ số nào là đáp án cần lấy.

\section{Ghi chú}

Bài viết đầy đủ có chứng minh và đề bài chi tiết. Bản slide hữu ích nhất ở
vai trò bảng nhắc: nhận dạng phép ghép tổ hợp nào tương ứng với phép toán nào
trên hàm sinh.

\end{document}
